\subsection{Remainders}
When we divide numbers, we sometimes get a remainder.
\begin{itemize}
\item With numbers, the remainder must be less than the divisor.
\item With polynomials, the \makeblank of the remainder must be less than the \makeblank of the divisor.
\end{itemize}
\begin{Example}
Consider the division problem:
\[
\left(2x^3-9x^2+14x-5\right)\div\left(x^2-3x+4\right)
\]
\end{Example}
The divisor has degree \makeblank[0.5in], so the remainder has degree at most \makeblank[0.5in].
\begin{lpic}{}{8}
  \divbracket{3}{4}{-2};
  \regtextright{3.25}{-3}{$x^2-3x+4$};
  \foreach \x/\lbl in {0/2x^3,2/{}-9x^2,4/{}+14x,6/{}-5}
    \regtextright{5+\x}{-3}{$\lbl$};
	\writethink[12]{-1}{8}
\end{lpic}
The remainder is \makeblank which has degree \makeblank[0.5in].

\noindent We write this result as:
\[
\left(2x^3-9x^2+14x-5\right)\div\left(x^2-3x+4\right)
=\makeblank[1.25in]+\makeblank[1.25in]
\]
This is just like a mixed number! Like we did with numbers, we can also write our result in terms of multiplication and addition.
\[
2x^3-9x^2+14x-5
=\Bigl(\makeblank[1in]\Bigr)
 \Bigl(\makeblank\Bigr)+
 \Bigl(\makeblank[1in]\Bigr)
\]
We can think of all these polynomials as functions, and write
\[
P(x)=Q(x)D(x)+R(x)
\]
\begin{itemize}
\item \makeblank[.5in] is the dividend, \makeblank
\item \makeblank[.5in] is the divisor, \makeblank
\item \makeblank[.5in] is the quotient, \makeblank
\item \makeblank[.5in] is the remainder, \makeblank
\end{itemize}
Remember, the degree of the \makeblank[1.5in] is always \makeblank[1in] than the degree of the \makeblank[1.5in].